\documentclass[11pt,a4paper]{article}

\usepackage{ctex}
\usepackage[T1]{fontenc}
\usepackage[top=2.5cm, bottom=2.5cm, left=2.5cm, right=2.5cm]{geometry}
\usepackage{booktabs}
\usepackage{array}
\usepackage{enumitem}
\usepackage{tabularx}
\usepackage{amssymb}
\usepackage{xcolor}
\definecolor{primary}{HTML}{1a1a2e}
\definecolor{accent}{HTML}{3b82f6}
\definecolor{success}{HTML}{22c55e}
\definecolor{danger}{HTML}{ef4444}
\definecolor{warning}{HTML}{f59e0b}
\definecolor{graytext}{HTML}{6b7280}
\definecolor{progress}{HTML}{8b5cf6}
\usepackage[hidelinks]{hyperref}
\hypersetup{colorlinks=true, linkcolor=primary, urlcolor=accent}
\usepackage{titlesec}
\titleformat{\section}{\Large\bfseries\color{primary}}{}{0em}{}[\titlerule]
\titleformat{\subsection}{\large\bfseries\color{primary}}{}{0em}{}
\usepackage{fancyhdr}
\pagestyle{fancy}
\fancyhf{}
\fancyhead[L]{\small\color{graytext} 企业级 AI Agent 应用商店 -- 需求文档}
\fancyhead[R]{\small\color{graytext} DEMAND-004}
\fancyfoot[C]{\small\color{graytext} \thepage}
\renewcommand{\headrulewidth}{0.4pt}
\usepackage{tikz}
\usetikzlibrary{shapes.geometric, arrows.meta, positioning, fit}

\title{\textbf{DEMAND-004：多智能体协作方案}\\[0.5em]\small OpenClaw / Claude / Codex / Cursor 的 Skill 适配}
\author{万能产品小助手（PM AI）\\魏子政 审阅}
\date{2026-05-06 / 版本 v0.1}

\begin{document}
\maketitle

\begin{center}
\begin{tabular}{>{\bfseries}p{3cm} p{10cm}}
\toprule
提出日期 & 2026-05-06 \\
提出者 & 魏子政 \\
优先级 & \textcolor{danger}{P0} \\
状态 & \textcolor{progress}{进行中} \\
总需求编号 & D-010 \\
\bottomrule
\end{tabular}
\end{center}

\vspace{1em}

\section{需求描述}

公司拥有多个代码工具和智能体（OpenClaw、Claude、Codex、Cursor 等）。需要设计一套多智能体协作方案，使各智能体能共享上下文、协同工作，而不是各自为战。

核心问题：每个智能体有自己的 Skill/规则加载机制，如何让它们在同一个项目中对齐认知？

\section{现状分析：各智能体的 Skill 机制}

\begin{table}[h]
\centering
\begin{tabularx}{\linewidth}{lXXXX}
\toprule
\textbf{智能体} & \textbf{Skill 加载方式} & \textbf{规则文件位置} & \textbf{触发机制} & \textbf{项目级支持} \\
\midrule
OpenClaw & SKILL.md（按需读取） & \texttt{skills/*/SKILL.md} & 对话匹配 + 手动 read & 支持 \\
Cursor & .mdc 规则文件 & \texttt{.cursor/rules/*.mdc} & alwaysApply / globs & 支持 \\
Claude (Code) & CLAUDE.md & \texttt{CLAUDE.md} & 自动加载 & 支持 \\
Codex (OpenAI) & system prompt & 无项目级文件 & 仅全局 prompt & 不支持 \\
\bottomrule
\end{tabularx}
\end{table}

\section{协作架构设计}

\subsection{核心思路：统一上下文源}

不依赖各智能体自身的 Skill 机制来共享信息，而是在项目目录中维护一套\textbf{统一的项目上下文文件}，各智能体通过各自能识别的入口文件读取。

\begin{center}
\begin{tikzpicture}[
  node distance=1cm,
  ctx/.style={rectangle, rounded corners, draw=success, fill=success!5, minimum width=3.5cm, minimum height=0.8cm, align=center, font=\small},
  agent/.style={rectangle, rounded corners, draw=primary, fill=primary!5, minimum width=2.5cm, minimum height=0.8cm, align=center, font=\small},
  arrow/.style={-{Stealth[length=2.5mm]}, thick, graytext}
]
  \node[ctx] (master) {MASTER.tex\\（总需求文档）};
  \node[ctx, below=of master] (guide) {GUIDE-001.tex\\（目录规范）};
  \node[ctx, below=of guide] (naming) {DOC-NAMING-CONVENTION.tex\\（命名规范）};

  \node[agent, right=2cm of master] (oc) {OpenClaw\\(PM AI)};
  \node[agent, right=2cm of guide] (cursor) {Cursor\\(开发)};
  \node[agent, right=2cm of naming] (claude) {Claude\\(辅助)};

  \draw[arrow] (master) -- (oc);
  \draw[arrow] (master) -- (cursor);
  \draw[arrow] (master) -- (claude);
  \draw[arrow] (guide) -- (cursor);
  \draw[arrow] (naming) -- (oc);
\end{tikzpicture}
\end{center}

\subsection{三层上下文体系}

\begin{table}[h]
\centering
\begin{tabularx}{\linewidth}{lp{3cm}X}
\toprule
\textbf{层} & \textbf{文件} & \textbf{说明} \\
\midrule
L1：全局认知 & CLAUDE.md / AGENTS.md & 技术栈锁定、项目概述、禁止项。所有智能体都能自动加载。 \\
L2：需求上下文 & MASTER.tex / DEMAND-xxx.tex & 当前需求全貌和历史。PM 需要时 read，Cursor 通过 mdc 引用。 \\
L3：执行规范 & GUIDE-xxx.tex / .mdc & 目录结构、编码规范、测试标准。Cursor 自动加载，Claude 参考 CLAUDE.md。 \\
\bottomrule
\end{tabularx}
\end{table}

\section{各智能体的适配策略}

\subsection{OpenClaw（PM AI）-- 协调者}

\begin{itemize}[nosep]
  \item \textbf{角色}：唯一的需求入口和协调中心。魏子政只与 OpenClaw 对话。
  \item \textbf{Skill 加载}：\texttt{skills/*/SKILL.md}（按需 read）
  \item \textbf{项目上下文}：通过 read 工具读取 MASTER.tex、DEMAND-xxx.tex、GUIDE-xxx.tex
  \item \textbf{调度其他智能体}：
    \begin{itemize}[nosep]
      \item 调度 Cursor：\texttt{cursor agent --print --trust --workspace <path> "任务"}
      \item 调度 Claude：通过 sessions\_spawn 或 API 调用
      \item 调度 Codex：通过 API 调用
    \end{itemize}
  \item \textbf{适配文件}：AGENTS.md（全局指令）、CLAUDE.md（共享给 Cursor/Claude）
\end{itemize}

\subsection{Cursor（开发 AI）-- 执行者}

\begin{itemize}[nosep]
  \item \textbf{角色}：主力开发，负责前后端代码实现。
  \item \textbf{规则加载}：\texttt{.cursor/rules/*.mdc}（alwaysApply + globs）
  \item \textbf{项目上下文}：CLAUDE.md + AGENTS.md 自动加载；通过 @ 引用可加载其他文件
  \item \textbf{适配文件}：
    \begin{itemize}[nosep]
      \item \texttt{CLAUDE.md}：技术栈、禁止项、mdc 索引
      \item \texttt{.cursor/rules/project-overview.mdc}：项目概述（alwaysApply）
      \item \texttt{.cursor/rules/backend-rules.mdc}：后端规范（globs: backend/**）
      \item \texttt{.cursor/rules/frontend-rules.mdc}：前端规范（globs: frontend/**）
      \item \texttt{.cursor/rules/testing.mdc}：自检清单（alwaysApply）
    \end{itemize}
\end{itemize}

\subsection{Claude（辅助 AI）-- 顾问}

\begin{itemize}[nosep]
  \item \textbf{角色}：辅助角色，适合做代码审查、架构咨询、复杂问题分析。
  \item \textbf{规则加载}：\texttt{CLAUDE.md}（自动加载）
  \item \textbf{项目上下文}：通过 CLAUDE.md 获取技术栈和项目概述；通过 OpenClaw 传递需求上下文
  \item \textbf{适配文件}：CLAUDE.md（与 Cursor 共享同一份）
  \item \textbf{调用方式}：OpenClaw 通过 sessions\_spawn 或 API 调用，传入需求上下文
\end{itemize}

\subsection{Codex（OpenAI）-- 备选}

\begin{itemize}[nosep]
  \item \textbf{角色}：备选执行者，适合快速原型或补丁任务。
  \item \textbf{限制}：不支持项目级规则文件，只能通过 system prompt 传入上下文。
  \item \textbf{适配方式}：OpenClaw 在调用时将关键技术约束作为 system prompt 的一部分传入。
\end{itemize}

\section{Skill 适配的关键原则}

\begin{enumerate}
  \item \textbf{单一事实源}：MASTER.tex 是需求的唯一权威来源，所有智能体的认知从这里派生。
  \item \textbf{CLAUDE.md 共享}：Cursor 和 Claude 共享同一份 CLAUDE.md，确保技术栈和约束一致。
  \item \textbf{OpenClaw 协调}：所有智能体间的信息传递通过 OpenClaw 中转，不直接通信。
  \item \textbf{上下文传递而非共享}：各智能体不共享会话上下文，OpenClaw 在调度时传递必要信息。
  \item \textbf{规则文件最小化}：每个智能体只维护它能自动加载的规则文件，避免重复。
\end{enumerate}

\section{落地步骤}

\begin{enumerate}
  \item 确保当前项目的 CLAUDE.md 内容完整，同时服务 Cursor 和 Claude
  \item 在 CLAUDE.md 中补充指向 MASTER.tex 和 GUIDE-xxx.tex 的引用说明
  \item 测试 OpenClaw 调度 Cursor 的完整链路（已通过 D-008 验证）
  \item 测试 OpenClaw 调度 Claude 的链路（待验证）
  \item 记录各智能体的 skill 加载效果，持续优化
\end{enumerate}

\section{验收标准}

\begin{itemize}
  \item[$\square$] 多智能体协作方案经魏子政确认
  \item[$\square$] CLAUDE.md 能同时服务 Cursor 和 Claude
  \item[$\square$] OpenClaw 调度 Cursor 链路稳定（已有验证）
  \item[$\square$] OpenClaw 调度 Claude 链路验证通过
  \item[$\square$] 各智能体对项目技术栈的认知一致
\end{itemize}

\section*{更新记录}

\begin{tabular}{lll}
\toprule
版本 & 日期 & 变更内容 \\
\midrule
v0.1 & 2026-05-06 & 初始版本 \\
\bottomrule
\end{tabular}

\end{document}
